\section{Introduction}
\label{sec:introduction}

A vulnerability identifier is a starting point for investigation, not an
executable experiment. To confirm a reported IoT vulnerability, an analyst
may have to locate an old firmware release, unpack a vendor-specific
image, identify the affected service, and reconstruct enough of the
device environment to deliver a meaningful test input. A working exploit
script is useful only when its behavior is connected to the intended
target. This distinction matters for validating vulnerability reports,
building regression tests, and deciding when automated security analysis
has produced actionable evidence.

LLM agents can search documentation, write programs, and invoke analysis
tools within one interactive workflow. Research has demonstrated their
potential in exploitation, penetration testing, and vulnerability
repair~\cite{fang2024oneday,happe2026enterprisenetworks,zhang2026ossfuzzrepair}.
However, the meaning of a reported success depends on what the experiment
provides. Testing an already running vulnerable application, constructing
a PoC from a patch, and locating and executing an affected firmware
component expose different combinations of information and execution
requirements. Their success rates answer different questions.

Existing work covers important parts of this space. CVE-Bench provides
real-world Web application targets and attack objectives; CyberGym
evaluates reproduction with prepared software environments; and
Patch-to-PoC studies kernel reproduction with patch input and dedicated
execution support~\cite{cvebench,cybergym,krepro}. Environment construction
is itself addressed by systems such as
CVE-GENIE~\cite{cvegenie}. Firmware analysis systems separately establish
that obtaining a binary does not remove its operating-system, peripheral,
and configuration dependencies~\cite{challenges,greenhouse,pandawan}.
These contributions motivate a closer empirical account of what general
coding agents do when they must navigate the IoT reproduction workflow
from a CVE identifier.

We study three connected questions. First, what can these agents actually
complete under the observed conditions? Second, where do their attempts
break down, and which unmet prerequisites are visible in the evidence?
Third, what do agents do after encountering an obstacle, and how do their
subsequent actions and reports relate to their actual progress? The third
question includes repairing a real execution environment, repeating an
unsuccessful strategy, changing the tested entry point, implementing a
simulated service, and stopping. Examining these responses can explain
differences between superficially similar final artifacts.

Consider two records from our archive. One agent extracts Xiaomi firmware
and sends a test input to an original service running in a chroot. Another
implements a Python server for a NETGEAR authentication-bypass pattern and
successfully tests that server. Both records contain executable scripts
and positive demonstration output. They nevertheless support different
claims about the specified firmware targets. Even the first record needs
a further distinction: reaching an internal component directly does not
establish that the documented external entry point and its authentication
conditions were reproduced. We use such distinctions to connect final
outcomes to the process that produced them.

Our material consists of 450 historical runs: 30 selected CVEs, five model
configurations, and three repetitions per combination. The experiments
used a shared OpenCode harness on an ARM64 host. Each agent received a
CVE identifier within a common task prompt and could obtain additional
public information and tools. We use six reproduction activities to
locate progress, while allowing iteration and alternative paths. Target
identity, observed execution, vulnerability-specific effects, and agent
reports are recorded separately. No weighted plan, task, or overall score
is used to answer the research questions.

\paragraph{Draft status.}
This manuscript is an internal working draft. Archive counts and explicitly
described artifact inspections are available; comprehensive outcome
validation and trajectory coding remain incomplete. Visible TODOs identify
unfinished analyses and experiments. Historical evaluator labels are
identified as such and are not presented as independently verified
reproduction results. \todo{INTRO-01}{Replace this draft-status paragraph
with a concise summary of the audited principal findings before submission.}

\paragraph{Contributions.}
The intended contributions are an empirical account of CVE-initiated IoT
reproduction under a documented agent setup; an evidence-based analysis
connecting outcomes, obstacles, and responses; and a research resource,
\sys{}, supporting inspection and reanalysis. The knowledge contribution
will be stated as concrete findings once the audits are complete.
\todo{INTRO-02}{Replace the intended-contribution wording with final,
specific findings and the actual scope of the released resource.}

Our conclusions concern the studied cases and agent configurations.
They do not estimate the prevalence of reproducible vulnerabilities in
all IoT products, isolate an intrinsic property of a model, or establish
the effect of removing an obstacle without an intervention. These
boundaries guide both the analysis and the practical implications.

